% c07 —— §8.2 周期区域 + §8.3 周期映射 + §8.4 开头（p017–p019 / 原书 p.143–145）
% 承 c06 末 (8.10)；本块从「且它在齐次坐标下可写成……」起，止于定理 8.2 第 2 条句完。

% ===== 原书 p.143 =====

\noindent 且它在齐次坐标下可写成
\begin{equation}
  \lambda^{2}\nu^{6}+4\mu\nu^{5}-12\lambda^{3}\nu^{4}-50\lambda\mu\nu^{3}
  +48\lambda^{4}\nu^{2}+1000\lambda^{2}\mu\nu-64\lambda^{5}+3125\mu^{2}. \tag{8.11}
\end{equation}
% [Tutor Note: 原书 PDF 文字层在 (8.11) 中对同一坐标混用了 $v$ 与 $\nu$（前三处作 v、后三处作 \nu；(8.22) 的第一个坐标亦作 v），本译本按 (8.9) 的记法统一记作 $\nu$。]
判别式 $\boldsymbol{\nabla}$ 的 Horn–Kapranov 单值化
$h\colon\mathbb{P}^{1}\to\mathbb{L}\otimes\mathbb{C}^{\times}$，$[\lambda_{1}:\lambda_{2}]\mapsto(\lambda,\mu)$，由
\begin{align}
  \lambda&=\frac{\lambda_{1}(\lambda_{1}+2\lambda_{2})}{(2\lambda_{1}+5\lambda_{2})^{2}}, \tag{8.12}\\
  \mu&=-\frac{\lambda_{2}^{3}(\lambda_{1}+2\lambda_{2})^{2}}{(2\lambda_{1}+5\lambda_{2})^{5}} \tag{8.13}
\end{align}
给出。

\subsection{周期区域}

设 $M$ 为 (8.1) 中给出的格，则 $M$ 极化 K3 曲面的模空间可按如下方式与对称 Hilbert 模曲面等同。令实二次域 $\mathbb{Q}(\sqrt{5})$ 的整数环为 $\mathcal{O}$。Hilbert 模群 $PSL_{2}(\mathcal{O})$ 按
\begin{equation}
  \begin{pmatrix}\alpha&\beta\\ \gamma&\delta\end{pmatrix}\colon
  (z_{1},z_{2})\mapsto
  \left(\frac{\alpha z_{1}+\beta}{\gamma z_{1}+\delta},\
        \frac{\alpha' z_{1}+\beta'}{\gamma' z_{1}+\delta'}\right), \tag{8.14}
\end{equation}
作用在两个上半平面的直积 $\mathbb{H}\times\mathbb{H}$ 上。这里 $(-)'$ 表示在 $\mathbb{Q}(\sqrt{5})$ 中的共轭。
% [Tutor Note: (8.14) 右端第二个分量原书印作 $z_{1}$（两个分量都写作 $z_{1}$）；按 $PSL_{2}(\mathcal{O})$ 在 $\mathbb{H}\times\mathbb{H}$ 上的作用，第二个分量应为 $z_{2}$，此处照原图录作 $z_{1}$。]

由于
\begin{equation}
  \begin{pmatrix}2&1\\1&-2\end{pmatrix}=WUW^{T},\qquad
  W=\begin{pmatrix}1&1\\ -\varepsilon^{-1}&\varepsilon\end{pmatrix},\qquad
  \varepsilon=(1+\sqrt{5})/2 \tag{8.15}
\end{equation}
所以
\begin{equation}
  \mathbb{H}\times\mathbb{H}\to\mathcal{D}^{+},\qquad
  (z_{1},z_{2})\mapsto(I_{2}\oplus(W^{T})^{-1})
  \begin{pmatrix}z_{1}z_{2}\\ -1\\ z_{1}\\ z_{2}\end{pmatrix} \tag{8.16}
\end{equation}
给出双全纯映射。在这一等同之下，格的正交群 $PO^{+}(T)$ 由 Hilbert 模群
$PSL_{2}(\mathcal{O})$ 与置换
\begin{equation}
  \tau\colon\mathbb{H}\times\mathbb{H}\to\mathbb{H}\times\mathbb{H},\qquad
  (z_{1},z_{2})\mapsto(z_{2},z_{1}) \tag{8.17}
\end{equation}
生成。对称 Hilbert 模曲面 $\mathbb{H}\times\mathbb{H}/\langle PSL_{2}(\mathcal{O}),\tau\rangle$
已由 Hirzebruch [Hir77] 详细研究（另见 [KKN89]）。对称 Hilbert 模形式构成的分次
% ===== 原书 p.144 =====
环 $\mathfrak{M}=\bigoplus_{n=0}^{\infty}\mathfrak{M}_{n}$ 由权为 2、6、10、15 的元素
$\mathfrak{A},\mathfrak{B},\mathfrak{C},\mathfrak{D}$ 生成，其关系式由权为 30 的式 (8.18)、(8.19) 生成。
\begin{equation}
  \mathfrak{M}=\mathbb{C}[\mathfrak{A},\mathfrak{B},\mathfrak{C},\mathfrak{D}]
  /(144\mathfrak{D}^{2}-\Delta(\mathfrak{A},\mathfrak{B},\mathfrak{C})), \tag{8.18}
\end{equation}
\begin{equation}
  \Delta(\mathfrak{A},\mathfrak{B},\mathfrak{C})
  =-1728\mathfrak{B}^{5}+720\mathfrak{A}\mathfrak{B}^{3}\mathfrak{C}
  -80\mathfrak{A}^{2}\mathfrak{B}\mathfrak{C}^{2}
  +64\mathfrak{A}^{3}(5\mathfrak{B}^{2}-\mathfrak{A}\mathfrak{C})^{2}
  +\mathfrak{C}^{3} \tag{8.19}
\end{equation}
因此，叠 $\overline{\mathcal{M}}=\Proj\mathfrak{M}=[(\Spec\mathfrak{M}\setminus 0)/\mathbb{C}^{\times}]$
是加权射影空间 $\mathbb{P}(2,6,10,15)$ 中的 30 次超曲面。该超曲面可从加权射影平面
$\mathbb{P}(1,3,5)=\Proj\mathbb{C}[\mathfrak{A},\mathfrak{B},\mathfrak{C}]$ 出发，沿由 $\Delta$
定义的除子作二阶根构造得到，其尖点由唯一一个点 $[\mathfrak{A}:\mathfrak{B}:\mathfrak{C}]=[1:0:0]$ 构成。

$M$ 极化 K3 曲面族的 Weierstrass 模型由
\begin{equation}
  z^{2}=x^{3}-4(4y^{3}-5\mathfrak{A}y^{2})x^{2}+20\mathfrak{B}xy^{3}+\mathfrak{C}y^{4} \tag{8.20}
\end{equation}
给出 [Nag13]。该族的两个纤维作为 $M$ 极化 K3 曲面同构的充要条件是：存在
$\alpha\in\mathbb{C}^{\times}$，使得
$(\mathfrak{A}',\mathfrak{B}',\mathfrak{C}')=(\alpha\mathfrak{A},\alpha^{3}\mathfrak{B},\alpha^{5}\mathfrak{C})$。

\subsection{周期映射}

若把周期映射 $X_{\mathcal{F}(A_{0})}\dashrightarrow\mathbb{P}(1,3,5)$ 写成有理映射
——即由二次叠 $X_{\mathcal{F}(A_{0})}$ 经加权收缩得到的 $\mathbb{P}(1,2,5)$ 到
$\mathcal{M}^{*}$ 的粗模空间 $\mathbb{P}(1,3,5)$ 的映射——
\begin{equation}
  \mathbb{P}(1,2,5)\dashrightarrow\mathbb{P}(1,3,5) \tag{8.21}
\end{equation}
则它由下式给出
\begin{equation}
  [\nu:\lambda:\mu]\mapsto[\mathfrak{A}:\mathfrak{B}:\mathfrak{C}]
  =\left[\lambda-\frac{\nu^{2}}{4}:\frac{25}{2}\nu\mu:-3125\mu^{2}\right] \tag{8.22}
\end{equation}
[Nag12, Theorem 6.2]。该有理映射以 $[\nu:\lambda:\mu]=[1:1/4:0]$ 为不定点。该不定点经沿理想
$(\lambda-\nu^{2}/4,\mu)$ 的权为 $(1,3)$ 的加权爆破而被消除，结果得到把因子
$\{\mu=0\}\subset\mathbb{P}(1,2,5)$ 的严格变换收缩为尖点 $[1:0:0]\in\mathbb{P}(1,3,5)$ 的权为
$(1,2)$ 的加权收缩。这一情形表示在图 8.3 中。

通过周期映射 (8.22)，由直接计算容易验证 (8.19) 与 (8.11) 一致。另一方面，加权射影平面
$\mathbb{P}(1,3,5)$ 的自同构由形如
\begin{equation}
  [x:y:z]\mapsto[ax:by+cx^{3}:dz+exy+fx^{5}] \tag{8.23}
\end{equation}
的加权射影变换组成，而其中把判别式簇映为判别式簇的只有恒等映射，这也容易看出。

% ===== 原书 p.145 =====

\begin{figure}[htbp]
  \centering
  \includegraphics[width=0.72\textwidth]{fig8_3.png}
  \caption{参数空间的双有理变换（原书 p.145）}
\end{figure}

\subsection{镜像对称与单值性}

$\Delta$ 的极多面体由
\[
  \check{\Delta}=\Conv\bigl\{(0,-1,1),(4,-1,-1),(-1,-1,-1),(-1,-1,1),((-1,4,-1),(-1,0,1)\bigr\}
\]
给出。
% [Tutor Note: 极多面体的顶点列中，原书在 $(-1,4,-1)$ 之前多印了一个左括号（印作「$((-1,4,-1),(-1,0,1)\}$」），括号不配对；此处照原图录出。]

环面簇 $\check{X}$ 是 $\mathbb{P}^{2}$ 上的 $\mathbb{P}^{1}$ 丛
$\mathbb{P}(\mathcal{O}_{\mathbb{P}^{2}}\oplus\mathcal{O}_{\mathbb{P}^{2}}(2))$，因而是一个 Fano 流形。
镜像族 $\check{\mathcal{M}}$ 的非常一般的纤维的 Picard 格，可由下面的 Noether–Lefschetz 型定理容易求出：

\begin{theorem}[{[Moĭ67, Theorem 7.5]}]
设 $V$ 为光滑三维射影簇，$\iota\colon E\hookrightarrow V$ 为非常一般的超平面截断。此时，
$\iota^{*}\colon\Pic(E)\to\Pic(V)$ 为满射的充要条件是，下列之一成立：
\begin{enumerate}[label=\arabic*.,leftmargin=2.2em]
  \item Betti 数满足 $b_{2}(V)=b_{2}(E)$。
  \item Hodge 数满足 $h^{2,0}(E)>h^{2,Q}(V)$。
\end{enumerate}
\end{theorem}
% [Tutor Note: 第 2 条右端的 Hodge 数原书印作 $h^{2,Q}(V)$：高倍核对可见上标为字母 $Q$（左端 $h^{2,0}(E)$ 的上标才是数字 $0$），PDF 文字层亦作 $h^{2,Q}$。按 Hodge 数记号应为 $h^{2,0}(V)$，此处照原图录作 $h^{2,Q}$，读者按 $h^{2,0}$ 理解即可。]

% CHUNK c07 DONE pages p017–p019
